\documentclass[11pt]{article}
\usepackage[margin=1in]{geometry}
\usepackage{amsmath,amssymb,amsthm}
\usepackage{graphicx}
\usepackage{hyperref}
\usepackage{microtype}
\usepackage{enumitem}
\usepackage{fancyvrb}

\newtheorem{theorem}{Theorem}
\newtheorem{lemma}{Lemma}
\newcommand{\CS}{\mathcal{CS}}
\newcommand{\Ucal}{\mathcal{U}}
\newcommand{\Ccal}{\mathcal{C}}
\newcommand{\T}{^{\mathsf T}}

\title{Pairwise-nonorthogonal complex symmetric matrices\\without truncated Toeplitz models}
\author{Candidate counterexample packet --- \texttt{agent\_lane\_04} (GPT5.6)}
\date{26 August 2026}

\begin{document}
\maketitle

\noindent\textbf{Status.}
Candidate counterexample, likely valid; human review needed.

\section{Source question}

Garcia, Poore, and Ross ask in Section~7, Question~2 of
\cite{GPR} (PDF page~13):
\begin{quote}
Let $n\geq4$. Is every complex symmetric matrix $M\in M_n(\mathbb C)$
having no pair of orthogonal, nonzero eigenvectors unitarily equivalent to a
truncated Toeplitz operator?
\end{quote}

\begin{center}
\includegraphics[width=0.97\textwidth]{figures/open_problem_crop.png}
\end{center}

Here ``complex symmetric matrix'' means $M\T=M$.  Orthogonality in the
question is for the usual Hermitian inner product on $\mathbb C^n$.

\section{Result}

\begin{theorem}\label{thm:main}
For every $n\geq10$, there exists $M\in M_n(\mathbb C)$ such that
\begin{enumerate}[label=(\roman*)]
\item $M\T=M$;
\item $M$ has simple spectrum and no two nonzero eigenvectors are orthogonal;
\item $M$ is not unitarily equivalent to a truncated Toeplitz operator.
\end{enumerate}
Consequently Question~2 of \cite{GPR} has a negative answer.
\end{theorem}

The theorem is existential.  It proves that counterexamples occupy a
nonempty part of a full-dimensional open semialgebraic locus after a
lower-dimensional exceptional set is removed; it does not identify a
particular rational or algebraic matrix outside that exceptional set.

\section{Proof intuition}

O'Loughlin associates to finite-dimensional truncated Toeplitz operators a
semialgebraic set $\Ucal(\Ccal)\subseteq\CS_n$ and proves that its real
dimension is strictly less than $n(n+1)$ for $n\geq10$ \cite{OL}.  It is not
enough merely to choose a generic irreducible symmetric matrix: the source
question imposes the sharper condition that every two eigenvectors be
nonorthogonal.

The key observation is that this sharper condition is still generic on the
simple-spectrum locus.  It is semialgebraic, open, and nonempty.  For
nonemptiness, start with a complex orthogonal matrix
$Q(t)=\exp(itK)$, where $K$ is real skew-symmetric with no zero
off-diagonal entries.  Its columns are orthogonal for the bilinear form
$x\T y$, but, for small $t\ne0$, every pair has nonzero Hermitian inner
product because
$Q(t)^*Q(t)=\exp(2itK)=I+2itK+O(t^2)$.  Giving those columns distinct
eigenvalues produces the required symmetric matrices.  A full-dimensional
open semialgebraic set cannot be contained in O'Loughlin's
lower-dimensional TTO-model locus.

\section{Proof}

Let
\[
\CS_n=\{M\in M_n(\mathbb C):M\T=M\},
\]
viewed as a real vector space of dimension $n(n+1)$.

Define $\mathcal G_n\subseteq\CS_n$ to consist of matrices with simple
spectrum for which eigenvectors belonging to distinct eigenvalues have
nonzero Hermitian inner product.  Because the spectrum is simple, this is
equivalent to saying that no pair of nonzero eigenvectors is orthogonal.

\begin{lemma}\label{lem:generic}
For every $n\geq2$, the set $\mathcal G_n$ is a nonempty open
semialgebraic subset of $\CS_n$.  In particular,
$\dim_{\mathbb R}\mathcal G_n=n(n+1)$.
\end{lemma}

\begin{proof}
The simple-spectrum locus is semialgebraic: it is defined by nonvanishing of
the discriminant of the characteristic polynomial.  Within that locus, the
failure of pairwise nonorthogonality is the existence of
$\lambda\ne\mu$ and unit vectors $x,y$ satisfying
\[
Mx=\lambda x,\qquad My=\mu y,\qquad x^*y=0.
\]
After real and imaginary parts are separated, these are polynomial equations
and the inequality $|\lambda-\mu|^2>0$.  Tarski--Seidenberg therefore shows
that the bad set, and hence $\mathcal G_n$, is semialgebraic.

On the simple-spectrum locus the eigenlines vary continuously locally.  The
finitely many Hermitian inner products between normalized eigenvectors thus
vary continuously up to harmless phase factors.  Their nonvanishing is an
open condition, so $\mathcal G_n$ is open in $\CS_n$.

It remains to show nonemptiness.  Choose a real skew-symmetric matrix
$K=(k_{ij})$ with $k_{ij}\ne0$ whenever $i\ne j$; for example take
$k_{ij}=1$ for $i<j$ and $k_{ij}=-1$ for $i>j$.  Put
\[
Q(t)=\exp(itK).
\]
Since $K\T=-K$, we have $Q(t)\T Q(t)=I$.  Moreover $itK$ is Hermitian,
so
\[
Q(t)^*Q(t)=Q(t)^2=\exp(2itK)=I+2itK+O(t^2).
\]
For each $i\ne j$, its $(i,j)$ entry has derivative $2ik_{ij}\ne0$ at
$t=0$.  Since there are only finitely many pairs, all off-diagonal entries
are nonzero for every sufficiently small positive $t$.

Fix such a $t$ and distinct complex numbers $\lambda_1,\ldots,\lambda_n$.
Set
\[
M=Q(t)\operatorname{diag}(\lambda_1,\ldots,\lambda_n)Q(t)\T.
\]
Then $M\T=M$, and the columns $q_j$ of $Q(t)$ satisfy
$Mq_j=\lambda_jq_j$, since $Q(t)\T Q(t)=I$.  Their Hermitian Gram matrix
is $Q(t)^*Q(t)$, whose off-diagonal entries are all nonzero.  Hence
$M\in\mathcal G_n$.

A nonempty open semialgebraic subset of the real vector space $\CS_n$ has
the full dimension $n(n+1)$, proving the last assertion.
\end{proof}

We now use the notation of O'Loughlin \cite[Section~4]{OL}.  Let
$\Ccal\subseteq\CS_n$ be the set of symmetric matrices $N$ for which there
is a diagonal matrix $D$ with distinct unit-circle entries and
$\operatorname{rank}[D,N]\leq2$, and let $\Ucal(\Ccal)$ be the set of
symmetric matrices unitarily equivalent to some element of $\Ccal$.

Theorem~3.1, Proposition~4.6, and the proof of Theorem~4.3 in \cite{OL}
give
\begin{equation}\label{eq:bound}
\dim_{\mathbb R}\Ucal(\Ccal)
\leq 7n-6+\frac{n(n-1)}2
<n(n+1),\qquad n\geq10.
\end{equation}
The same paper's Theorem~4.1 and Lemma~4.2 show that every finite-dimensional
truncated Toeplitz operator has a symmetric matrix representative in
$\Ccal$.  Consequently, any symmetric matrix unitarily equivalent to a
truncated Toeplitz operator belongs to $\Ucal(\Ccal)$.

For $n\geq10$, Lemma~\ref{lem:generic} and \eqref{eq:bound} imply
\[
\mathcal G_n\not\subseteq\Ucal(\Ccal),
\]
because a semialgebraic subset cannot contain another semialgebraic set of
strictly larger dimension.  Choose
$M\in\mathcal G_n\setminus\Ucal(\Ccal)$.  It is symmetric, has simple
spectrum and pairwise nonorthogonal eigenvectors, and cannot be unitarily
equivalent to a truncated Toeplitz operator.  This proves
Theorem~\ref{thm:main}.

\section{Verification}

\paragraph{Logical checks.}
The dimension comparison is strict already at $n=10$:
\[
7(10)-6+\frac{10\cdot9}{2}=109<110=10\cdot11.
\]
The difference grows thereafter.  The TTO implication uses only the
supporting paper's stated matrix representation and rank-two commutator
condition; no converse is assumed.

The semialgebraic description quantifies over normalized eigenvectors, so
zero vectors cannot enter the projected bad set.  Simple spectrum also
handles two eigenvectors for the same eigenvalue: they are scalar multiples
and cannot be orthogonal unless one is zero.

\paragraph{Computational sanity check.}
The included script constructs the displayed $Q(t)$ and $M$ for $n=10$,
$t=0.01$, and $\lambda_j=j$.  Run from the repository root with
\begin{Verbatim}[fontsize=\small]
conda run --no-capture-output -n sandbox python \
  runs/fa_banach_001/solutions/counterexamples/\
1012.4820_pairwise_nonorthogonal_tto_counterexamples/\
code/verify_open_locus_example.py
\end{Verbatim}
It verifies complex orthogonality, symmetry, eigenvector residuals,
nonvanishing Hermitian cross-inner-products, and the dimension arithmetic.
The recorded output is:
\VerbatimInput[fontsize=\small]{verification_output.txt}
This computation proves neither the semialgebraic dimension theorem nor that
this particular numerical matrix lies outside $\Ucal(\Ccal)$; it is only a
sanity check for Lemma~\ref{lem:generic}.

\section{Novelty check and limitations}

On 26 August 2026, the four lightweight run indexes were searched for
arXiv:1012.4820, the exact Question~2 wording, ``pairwise nonorthogonal
eigenvectors,'' and truncated Toeplitz terminology.  A bounded arXiv/web
search used the exact phrases ``no pair of orthogonal'' and ``pairwise
nonorthogonal eigenvectors,'' together with arXiv:1012.4820 and
arXiv:2607.14019.  It found the source paper and O'Loughlin's later dimension
paper, but no source stating the strengthened conclusion of
Theorem~\ref{thm:main}.  Novelty is therefore provisional rather than
certified.

The argument is existential.  It answers the universal question negatively
and supplies counterexamples in every dimension $n\geq10$, but it does not
settle dimensions $4$ through $9$.  It does not address Question~4 in
\cite{GPR}, concerning matrices $A\oplus A\T$.

\paragraph{Human-review recommendation.}
Send to a reviewer familiar with semialgebraic dimension and truncated
Toeplitz matrix models.  The two main audit points are: (1) the projected
semialgebraic description and openness of $\mathcal G_n$; and (2) the precise
implication from a TTO model to membership in $\Ucal(\Ccal)$.

\begin{thebibliography}{9}
\bibitem{GPR}
S. R. Garcia, D. E. Poore, and W. T. Ross,
\emph{Unitary equivalence to a truncated Toeplitz operator: analytic symbols},
Proc. Amer. Math. Soc. 140 (2012), 1281--1295;
arXiv:\href{https://arxiv.org/abs/1012.4820}{1012.4820}.

\bibitem{OL}
R. O'Loughlin,
\emph{Semialgebraic Dimension and Truncated Toeplitz Models for Complex
Symmetric Matrices},
arXiv:\href{https://arxiv.org/abs/2607.14019}{2607.14019}, 2026.
\end{thebibliography}

\end{document}
